\documentclass[11pt]{article}

\usepackage[a4paper,margin=25mm,headheight=15pt]{geometry}
\usepackage[T1]{fontenc}
\usepackage[utf8]{inputenc}
\usepackage{lmodern}
\usepackage{microtype}
\usepackage{amsmath}
\usepackage{booktabs}
\usepackage{array}
\usepackage{longtable}
\usepackage{enumitem}
\usepackage{xcolor}
\usepackage{fancyhdr}
\usepackage{tikz}
\usetikzlibrary{arrows.meta,positioning,shapes.geometric}
\usepackage[hidelinks]{hyperref}

\definecolor{Accent}{HTML}{315B7D}
\definecolor{NoteBlue}{HTML}{EAF2F8}
\definecolor{NoteBorder}{HTML}{7BA5C4}
\definecolor{Muted}{HTML}{5B6570}

\hypersetup{
  pdftitle={Materials and Methods (Module A, _PK): ancestry sorting and genomic architecture},
  pdfauthor={Poikela et al.},
  pdfsubject={Hybrid genome evolution in Formica wood ants}
}

\pagestyle{fancy}
\fancyhf{}
\lhead{\small\color{Muted}Materials and Methods (Module A)}
\rhead{\small\color{Muted}Poikela et al.}
\cfoot{\small\thepage}
\renewcommand{\headrulewidth}{0.4pt}

\setlength{\parindent}{0pt}
\setlength{\parskip}{0.65em}
\setlist[itemize]{leftmargin=1.5em,itemsep=0.2em,topsep=0.25em}
\setlist[enumerate]{leftmargin=1.7em,itemsep=0.25em,topsep=0.25em}
\renewcommand{\arraystretch}{1.15}
\renewcommand{\thesection}{S\arabic{section}}
\renewcommand{\thesubsection}{\thesection.\arabic{subsection}}

\newcommand{\Faq}{\textit{F.~aquilonia}}
\newcommand{\Fpo}{\textit{F.~polyctena}}
\newcommand{\rsq}{\ensuremath{r^{2}}}
\newcommand{\todo}[1]{%
  \par\noindent
  \fcolorbox{NoteBorder}{NoteBlue}{%
    \parbox{\dimexpr\linewidth-2\fboxsep-2\fboxrule\relax}{%
      \textbf{Open item.} #1}}%
  \par}

\title{\vspace{-3em}
  \textbf{\color{Accent}Materials and Methods}\\[0.3em]
  \Large Ancestry sorting and genomic architecture}
\author{Poikela et al.}
\date{\small Module A (\texttt{\_PK} pipeline)}

\begin{document}
\maketitle
\vspace{-1.5em}

\begin{center}
\small\color{Muted}
Working draft for the \texttt{\_PK} pipeline. It covers the extent, direction and
genomic architecture of ancestry sorting, using the LD-reduced units produced in
Module~0 (LD decay and complexity reduction). Climate association and
dependence-aware (block-bootstrap) uncertainty are documented separately. Blue
boxes mark items that remain open.
\end{center}

\section{Scope and analytical principles}

These methods describe the quantification of the extent and direction of ancestry
sorting and the analysis of its genomic architecture, using the LD-reduced units
from Module~0 (LD decay and complexity reduction).

Three principles guided the analyses. First, markers in LD carry partly
redundant information. We therefore used clusters of strongly correlated
markers as LD-reduced analytical units wherever possible, while retaining
marker-level analyses as diagnostics of spatial pseudoreplication. We do not
assume that the resulting units are strictly statistically independent.
Second, a variable evaluated as a predictor was not also used to select the
analysed data. In particular, loci were gated by parental allele frequency
rather than by the diagnostic index (DI), leaving the full range of DI
available for analysis. Third, these analyses do not assume complete demographic
independence of the twenty hybrid populations; analyses that use variation among
populations are interpreted as descriptive associations.

\todo{The general sample description currently reports 165 sequenced hybrid
individuals, whereas the analysis objects contain 164. Confirm which individual
was excluded, most likely because of excessive missing data, and report the
identifier and exclusion criterion in the final text. Until then, 164 is used
for all analysis-specific counts.}

\section{Extent and direction of ancestry sorting}

\subsection{Near-fixation and ancestry orientation}

For each marker or cluster consensus, dosage was oriented as the
\Faq-allele frequency using the sign of the parental allele-frequency
difference. Units for which the parents did not differ had no defined parental
direction. Within each hybrid population, a unit was classified as near-fixed
toward \Faq{} when the oriented allele frequency was at least \(\phi\),
near-fixed toward \Fpo{} when it was at most \(1-\phi\), and otherwise not
near-fixed. Here \(\phi=0.90\) is the frequency the fixed-for (major) allele must
reach, equivalently the requirement that the opposing allele fall to at most
\(1-\phi=0.10\). The same threshold was used throughout because the attainable
frequency values depend on the number of sampled diploid individuals in a
population.

\subsection{Sorting decomposition and classification}

Let \(n_{\mathrm{obs}}\) be the number of populations with an observed and
oriented frequency, and \(n_{\mathrm{aqu}}\) and \(n_{\mathrm{pol}}\) the
numbers near-fixed toward \Faq{} and \Fpo{}, respectively. The total proportion
near-fixed was decomposed as
\begin{equation}
\frac{n_{\mathrm{aqu}}+n_{\mathrm{pol}}}{n_{\mathrm{obs}}}
=
\left|\frac{n_{\mathrm{aqu}}-n_{\mathrm{pol}}}{n_{\mathrm{obs}}}\right|
+
\frac{2\min(n_{\mathrm{aqu}},n_{\mathrm{pol}})}{n_{\mathrm{obs}}},
\label{eq:sorting}
\end{equation}
where the first term is the unidirectional component and the second the
bidirectional component. With \(\tau=0.5\), a unit was classified as
unidirectional when the absolute unidirectional component was at least the
bidirectional component and at least \(\tau\); bidirectional when the
bidirectional component exceeded the absolute unidirectional component and was
at least \(\tau\); and unsorted otherwise. Exchanging the parental labels
changes the direction but not the magnitude or class.

The two thresholds were treated separately. The near-fixation floor \(\phi\) is
the per-population convention and was held constant; the sorting threshold
\(\tau\) is the cross-population comparison. Their influence on the results was
evaluated over a grid of \(\tau\in\{0.5,0.6,0.7,0.8\}\) and
\(\phi\in\{0.80,0.85,0.90,0.95\}\): the estimated prevalence of sorting declines
monotonically with either threshold, but the composition (predominantly
unidirectional) and the diagnostic-index direction reversal are preserved
throughout (threshold-sensitivity supplement, Figs~S1--S3).

Nouhaud et al.~\cite{nouhaud2022} evaluated predictable ancestry sorting by
asking whether the same ancestry component repeatedly fixed across three hybrid
populations and summarising the outcome in fixed 20-kb windows. The present
classification retains that biological contrast but extends it to twenty
populations, estimates the outcome directly from population allele frequencies,
and separates consistent-direction from opposing-direction near-fixation.
Applying the statistic first to SNPs and then to LD-reduced units also avoids
assigning equal analytical status to physical windows that may contain very
different amounts of redundant information.

\subsection{Parental polymorphism gate}

Near-fixation was interpreted as sorting only for variants segregating in the
founding parental pool. Analyses were therefore restricted to units with pooled
parental MAF \(\geq0.15\), calculated from pooled parental allele counts and
folded to the minor allele. This gate excludes variants already close to
monomorphic in the parental pool without selecting on DI, which is evaluated as
a predictor in subsequent analyses.

\subsection{Cluster-level analysis and dilution check}

The sorting statistic was first calculated per marker. It was then recalculated
for eMLG consensus genotypes so that a cluster contributed one aggregate
observation. For clusters without a stored consensus but containing a
marker-level sorting signal, a consensus was constructed by the same procedure
used for larger clusters. Parental consensus genotypes reused the reference
marker and allele-flip decisions derived from the hybrid genotypes, preventing
hybrid and parental consensuses from being placed on inconsistent orientations.
Cluster DI was aggregated as the maximum DI among member markers.

The cluster-level aggregation analysis is conditional: clusters entered this
comparison because they contained at least one marker with a SNP-level sorting
signal. Consequently, the proportion of aggregated clusters classified as
unsorted does not estimate the genome-wide fraction of unsorted clusters.
Instead, it asks whether an apparent marker-level sorting signal remains after
correlated members are represented jointly.

When a cluster containing an individually sorted marker was unsorted after
aggregation, consensus fidelity was used to distinguish lack of an aggregate
signal from possible dilution in a heterogeneous consensus. Low-fidelity
clusters were flagged rather than automatically interpreted as evidence against
sorting.

\subsection{Models of sorting direction}

Among unidirectionally sorted units, the probability of sorting toward \Faq{}
was modelled by logistic regression. The genome-wide direction model is
described with the architecture analyses below. The robustness of the
diagnostic-index direction reversal to the classification thresholds is shown in
the threshold-sensitivity supplement (Figs~S2 and~S3). Associations are
interpreted descriptively and are not assigned a causal mechanism from these
regressions alone.

\section{Genomic architecture of sorting}

\subsection{Marker-level genomic summaries}

Map-based recombination rates were assigned to markers by linear interpolation and expressed in cM/Mb. For each retained marker, let \(p_{\mathrm{aqu}}\) and \(p_{\mathrm{pol}}\) denote the allele frequencies in the \Faq{} and \Fpo{} reference samples, respectively, and let
\[
\bar p=\frac{p_{\mathrm{aqu}}+p_{\mathrm{pol}}}{2}.
\]
We calculated the mean expected heterozygosity within the parental species as
\[
H_S=
\frac{
2p_{\mathrm{aqu}}(1-p_{\mathrm{aqu}})
+
2p_{\mathrm{pol}}(1-p_{\mathrm{pol}})
}{2},
\]
and expected heterozygosity based on the equally weighted pooled parental allele frequency as
\[
H_T=2\bar p(1-\bar p).
\]
The between-parent allelic-difference summary was calculated as
\[
d_{xy}
=
p_{\mathrm{aqu}}(1-p_{\mathrm{pol}})
+
p_{\mathrm{pol}}(1-p_{\mathrm{aqu}}),
\]
and relative parental differentiation as
\[
F_{\mathrm{ST}}=\frac{H_T-H_S}{H_T},
\]
for markers with \(H_T>0\).

These quantities are marker-level summaries conditional on the polymorphic sites retained for analysis. Because invariant sites were not included, \(H_S\) and \(d_{xy}\) should not be interpreted as genome-wide nucleotide diversity or absolute sequence divergence, respectively.

\subsection{Marker-level and LD-reduced summaries}

Sorting was summarised across recombination-rate deciles twice: once using one
representative per LD-reduced unit and once using a random sample of individual
markers. This comparison was used to diagnose spatial pseudoreplication.
Differences between the two summaries indicate that a marker-level pattern is
influenced by unequal redundancy across the recombination landscape.

\subsection{Association models}

Relationships among DI, recombination rate, parental expected heterozygosity,
between-parent allelic difference, relative differentiation, pooled parental
MAF, and cluster size were first summarised using rank correlations and
recombination deciles (architecture results summary, Fig.~1a).

The magnitude of sorting was modelled against recombination rate at the
LD-reduced-unit level. Cluster size was not included in this model because it
can be a consequence of both recombination and local marker variability;
conditioning on it could induce an association between its causes.

Among unidirectionally sorted units, sorting direction was analysed by logistic
regression with standardised DI, standardised \(\log_{10}(\mathrm{recombination}
+0.1)\), standardised pooled parental MAF, and standardised log cluster size as
predictors (architecture results summary, Fig.~1c). The recombination
coefficient is treated as an association rather than as evidence of linked
selection; its stability across the sorting threshold is shown in the
threshold-sensitivity supplement (Fig.~S3).

\subsection{Consensus-versus-representative validation}

For clusters represented by both an eMLG consensus and a representative marker,
we compared the inferred sorting direction between representations. The two
representations agreed on the parental direction for approximately 99\% of these
clusters, confirming that using representatives for genome-wide architecture models
does not materially change the parental direction assigned to clusters, while the
consensus remains preferable for aggregate sorting magnitude and cluster-level LD
analyses.

\section{Statistical interpretation}

The analyses are descriptive and conditional on the sampled populations,
quality gates, and analysis definitions. Consistent sorting can be produced by
more than one evolutionary process; distinguishing drift from selection requires
the recombination-matched neutral null, because a non-recombining autosomal
block does not have lower \(N_e\) than its surroundings under neutrality. The
weak residual low-recombination pull toward \Faq{} that survives control for DI
is therefore reported as a lead for the intrinsic hypothesis rather than as
evidence of linked selection.

\section{Software, reproducibility, and current status}

Ancestry-sorting and architecture analyses were performed in R using scripts in
the \texttt{formica\_hybrid} repository, folder \texttt{moduleA\_sorting/R/}:
\texttt{moduleA\_sorting\_phenomenon.R} (extent and direction of sorting, eMLG
consensus construction, dilution check, and the \(\tau\)/\(\phi\)
threshold-sensitivity sweeps) and \texttt{moduleA\_architecture.R} (genomic
architecture and sorting-direction models); shared statistical functions are in
\texttt{Ohta.R}, \texttt{parallelism\_stats.R} and \texttt{eMLG\_parallelism.R}.
Shared genotype inputs and the Module~0 clustering are read from \texttt{data/}
and \texttt{module0\_ld\_pruning/data/}; module intermediates are cached to
\texttt{moduleA\_sorting/data/} with a parameter-stamped validity check so that
reruns reuse them when settings are unchanged; figures are written to
\texttt{moduleA\_sorting/Figures/}.

\todo{Add the final \texttt{LDscnR} citation, version, and repository or archive
identifier before submission.}

\begin{longtable}{>{\raggedright\arraybackslash}p{0.20\textwidth}
                  >{\raggedright\arraybackslash}p{0.42\textwidth}
                  >{\raggedright\arraybackslash}p{0.27\textwidth}}
\caption{Principal parameter settings used in the Module A analyses.}
\label{tab:parameters}\\
\toprule
Analysis & Parameter & Value \\
\midrule
\endfirsthead
\toprule
Analysis & Parameter & Value \\
\midrule
\endhead
Sorting & Near-fixation floor & \(\phi=0.90\) (fixed-for allele) \\
        & Classification threshold & \(\tau=0.5\) \\
        & Parental MAF gate & \(\geq0.15\) \\
        & Cluster DI aggregation & maximum over members \\
\midrule
Threshold & Sorting threshold grid & \(\tau\in\{0.5,0.6,0.7,0.8\}\) \\
sensitivity & Near-fixation floor grid & \(\phi\in\{0.80,0.85,0.90,0.95\}\) \\
\midrule
Architecture & DI criterion (diagnostic) & \(\mathrm{DI}>-25\) \\
             & Direction model & logistic: DI, \(\log_{10}\)(recomb+0.1), parental MAF, \(\log\) size \\
             & Magnitude model & \(\mathrm{prop\_fixed}\sim\) recombination (+ DI) \\
\bottomrule
\end{longtable}

\begin{thebibliography}{1}
\bibitem{nouhaud2022}
Nouhaud P, Martin SH, Portinha B, Sousa VC, Kulmuni J. Rapid and
predictable genome evolution across three hybrid ant populations.
\textit{PLoS Biology}. 2022;20:e3001914.
\href{https://doi.org/10.1371/journal.pbio.3001914}
{doi:10.1371/journal.pbio.3001914}.
\end{thebibliography}

\end{document}
